% =====================================================================
% Parte 09 — Desigualdades (Stewart, Precálculo 7.ª ed., §1.8, pp. 81–91)
%
% Ejercicios elegidos (número del libro — página del libro):
%   15  (p. 89)  — lineal simple
%   33  (p. 89)  — desigualdad compuesta (doble)
%   41  (p. 89)  — cuadrática por tabla de signos
%   43  (p. 89)  — cuadrática con coeficiente principal 2
%   55  (p. 89)  — polinomial con factor repetido
%   57  (p. 89)  — cúbica x^3 - 4x > 0
%   59  (p. 89)  — racional ya factorizada
%   61  (p. 89)  — racional con constante al otro lado (trampa clásica)
%   63  (p. 89)  — racional con constante al otro lado
%   79  (p. 89)  — |x - 5| <= 3
%   83  (p. 89)  — |3x - 2| >= 5
%   89  (p. 89)  — 8 - |2x - 1| >= 6
%   93  (p. 89)  — frase --> desigualdad con valor absoluto
%  101  (p. 89)  — dominio de raíz: x^2 - 9 >= 0
%  111  (p. 90)  — aplicación: costo de la renta de un auto
%
% Sub-bloques cubiertos: 13–36 lineales, 37–58 no lineales, 59–74 cocientes,
% 75–90 valor absoluto, 91–94 frase --> valor absoluto, 101–104 dominio,
% 109–120 aplicaciones.
%
% Figuras extraídas del PDF: ninguna. Las rectas numéricas se dibujan con TikZ.
% =====================================================================

\documentclass[11pt]{article}
\usepackage[utf8]{inputenc}
\usepackage[spanish]{babel}
\usepackage[margin=2.3cm]{geometry}
\usepackage{amsmath,amssymb}
\usepackage{enumitem}
\usepackage{tcolorbox}
\tcbuselibrary{skins,breakable}
\usepackage{pgfplots}
\pgfplotsset{compat=1.18}
\usepackage{tikz}
\usepackage{graphicx}
\usepackage{array}
\usepackage[hidelinks]{hyperref}

% Notación española opcional (descomenta si quieres \sen en lugar de \sin)
% \DeclareMathOperator{\sen}{sen}
% \DeclareMathOperator{\tg}{tg}
% \DeclareMathOperator{\ctg}{ctg}

\newtcolorbox{sol}{colback=blue!5,colframe=blue!50!black,title=Solución paso a paso,breakable}
\newtcolorbox{teo}{colback=orange!8,colframe=orange!70!black,title=Teoría,breakable}
\newtcolorbox{nota}{colback=green!5,colframe=green!50!black,title=Idea clave,breakable}
\newtcolorbox{tip}{colback=yellow!10,colframe=yellow!60!black,title=Consejo,breakable}
\newtcolorbox{resp}{colback=red!5,colframe=red!50!black,title=Respuesta,breakable}
\newtcolorbox{reg}{colback=red!5,colframe=red!60!black,title=Regla,breakable}
\newtcolorbox{paso}[1]{colback=gray!5,colframe=gray!50!black,title={#1},breakable}

\graphicspath{{../images/}{images/}}
\newcommand{\sen}{\operatorname{sen}}

% babel spanish activa < y > y rompe TikZ/pgfplots (->, >=stealth)
\AtBeginDocument{\shorthandoff{<>}}

\begin{document}

\section*{Parte 09 --- Desigualdades (Stewart §1.8, pp. 81--91)}

\begin{teo}
\textbf{Guía del libro (p.~83) para desigualdades no lineales.}
\begin{enumerate}[leftmargin=*]
\item \textbf{Mueva todos los términos a un lado.} Reescriba la desigualdad con $0$ de un lado; si hay cocientes, únalos con común denominador. La técnica de factorización solo funciona en esta forma.
\item \textbf{Factorice} el lado distinto de cero.
\item \textbf{Encuentre los intervalos.} Los valores frontera son los ceros de cada factor (numerador \emph{y} denominador); ellos parten $\mathbb{R}$ en intervalos.
\item \textbf{Haga la tabla de signos} con un \emph{valor de prueba} por intervalo: cada factor conserva su signo dentro de un intervalo porque solo cambia de signo en su propio cero.
\item \textbf{Resuelva:} lea el último renglón y verifique aparte los puntos extremos (pueden entrar si el símbolo es $\le$ o $\ge$, y nunca entran si anulan un denominador).
\end{enumerate}
\end{teo}

\begin{tip}
Herramientas que se repiten en toda la parte.
\begin{itemize}[leftmargin=*]
\item \textbf{Reglas para desigualdades (p.~82):} sumar o restar lo mismo en ambos lados conserva el sentido; multiplicar o dividir por un número \emph{positivo} lo conserva; por un número \emph{negativo} lo \emph{invierte}.
\item \textbf{Signo de un producto o cociente (p.~83):} un número \emph{par} de factores negativos da valor positivo; un número \emph{impar}, negativo.
\item \textbf{Propiedades del valor absoluto (p.~86), con $c>0$:}
\[
|x|<c \iff -c<x<c, \qquad |x|\le c \iff -c\le x\le c,
\]
\[
|x|>c \iff x<-c \ \text{ o } \ c<x, \qquad |x|\ge c \iff x\le -c \ \text{ o } \ c\le x.
\]
Valen con cualquier expresión algebraica en lugar de $x$.
\item \textbf{Nunca} multiplique una desigualdad racional por su denominador: se desconoce su signo y no se sabe si hay que invertir el sentido (advertencia del libro, p.~85). Se usa tabla de signos.
\item \textbf{Notación de intervalos:} extremo incluido $\to$ corchete; extremo excluido o infinito $\to$ paréntesis.
\end{itemize}
\end{tip}

% =====================================================================
\subsection*{Ejercicio 1 --- Desigualdad lineal (libro 15, p. 89)}

\begin{quote}
\textbf{13--36 $\blacksquare$ Desigualdades lineales} \; Resuelva la desigualdad lineal. Exprese la solución usando notación de intervalos y trace la gráfica del conjunto solución.

\textbf{15.} $2x-5>3$
\end{quote}

\begin{paso}{Paso 1 --- Aislar el término con $x$}
\textbf{Qué hago:} sumo $5$ a ambos lados.

\textbf{¿De dónde sale?} El $-5$ es el término independiente del lado izquierdo del enunciado.

\textbf{¿Por qué?} Sumar la misma cantidad a ambos lados produce una desigualdad equivalente (regla 1, p.~82).
\[
2x-5>3 \;\Longleftrightarrow\; 2x>8.
\]
\end{paso}

\begin{paso}{Paso 2 --- Despejar $x$}
\textbf{Qué hago:} divido entre $2$ y conservo el sentido de la desigualdad.

\textbf{¿De dónde sale?} $2$ es el coeficiente de $x$ en $2x>8$.

\textbf{¿Por qué?} Dividir entre un número positivo no invierte el sentido (regla 3, p.~82).
\[
x>4.
\]
\end{paso}

\begin{center}
\begin{tikzpicture}[scale=0.75]
\draw[->] (-1,0) -- (8,0);
\foreach \x in {0,2,4,6} \draw (\x,0.12)--(\x,-0.12) node[below]{$\x$};
\draw[very thick,red] (4,0) -- (7.6,0);
\draw[red,fill=white,thick] (4,0) circle (2.6pt);
\end{tikzpicture}
\end{center}

\begin{resp}
$x>4$, es decir $(4,\infty)$.
\end{resp}

% =====================================================================
\subsection*{Ejercicio 2 --- Desigualdad compuesta con coeficiente negativo (libro 33, p. 89)}

\begin{quote}
\textbf{33.} $-2<8-2x\le -1$
\end{quote}

\begin{paso}{Paso 1 --- Restar $8$ en los tres miembros}
\textbf{Qué hago:} resto $8$ simultáneamente en los tres lados para dejar solo $-2x$ en el centro.

\textbf{¿De dónde sale?} El $8$ es el término constante del miembro central $8-2x$.

\textbf{¿Por qué?} La desigualdad doble equivale a las dos desigualdades a la vez, y restar lo mismo las conserva (regla 1, p.~82; ejemplo 2, p.~83).
\[
-2-8<-2x\le -1-8 \;\Longleftrightarrow\; -10<-2x\le -9.
\]
\end{paso}

\begin{paso}{Paso 2 --- Dividir entre $-2$ e invertir}
\textbf{Qué hago:} divido los tres miembros entre $-2$ e invierto ambos símbolos.

\textbf{¿De dónde sale?} $-2$ es el coeficiente de $x$ en $-2x$.

\textbf{¿Por qué?} Multiplicar o dividir por un número negativo invierte el sentido de la desigualdad (regla 4, p.~82).
\[
5>x\ge \tfrac92 \;\Longleftrightarrow\; \tfrac92\le x<5.
\]
\end{paso}

\begin{paso}{Paso 3 --- Verificación de un punto interior}
\textbf{Qué hago:} pruebo $x=4.6$, que está entre $4.5$ y $5$.

\textbf{¿De dónde sale?} Del intervalo obtenido en el paso 2.

\textbf{¿Por qué?} Sustituir en el enunciado original detecta un error de signo al invertir.
\[
8-2(4.6)=8-9.2=-1.2, \qquad -2<-1.2\le -1 \ \checkmark
\]
\end{paso}

\begin{center}
\begin{tikzpicture}[scale=0.75]
\draw[->] (-1,0) -- (7,0);
\foreach \x in {0,2,4,6} \draw (\x,0.12)--(\x,-0.12) node[below]{$\x$};
\draw (4.5,0.12)--(4.5,-0.12) node[below]{$\frac92$};
\draw[very thick,red] (4.5,0) -- (5,0);
\fill[red] (4.5,0) circle (2.6pt);
\draw[red,fill=white,thick] (5,0) circle (2.6pt);
\end{tikzpicture}
\end{center}

\begin{resp}
$\left[\dfrac92,\,5\right)$.
\end{resp}

% =====================================================================
\subsection*{Ejercicio 3 --- Cuadrática por tabla de signos (libro 41, p. 89)}

\begin{quote}
\textbf{37--58 $\blacksquare$ Desigualdad no lineal} \; Resuelva la desigualdad no lineal. Exprese la solución usando notación de intervalos y trace la gráfica del conjunto solución.

\textbf{41.} $x^2-3x-18\le 0$
\end{quote}

\begin{paso}{Paso 1 --- Factorizar el lado izquierdo}
\textbf{Qué hago:} factorizo el trinomio buscando dos números de producto $-18$ y suma $-3$: son $-6$ y $3$.

\textbf{¿De dónde sale?} Los coeficientes $-3$ y $-18$ del enunciado; todos los términos ya están de un lado con $0$ del otro (paso 1 de la guía, ya cumplido).

\textbf{¿Por qué?} La guía de la p.~83 exige factorizar para poder analizar el signo factor por factor.
\[
x^2-3x-18=(x-6)(x+3)\le 0.
\]
\end{paso}

\begin{paso}{Paso 2 --- Valores frontera e intervalos}
\textbf{Qué hago:} igualo cada factor a cero: $x=6$ y $x=-3$.

\textbf{¿De dónde sale?} De los factores $x-6$ y $x+3$ del paso 1.

\textbf{¿Por qué?} Un factor solo cambia de signo en su propio cero, así que esos puntos parten la recta en intervalos de signo constante.
\[
(-\infty,-3),\quad (-3,6),\quad (6,\infty).
\]
\end{paso}

\begin{paso}{Paso 3 --- Tabla de signos con valores de prueba}
\textbf{Qué hago:} evalúo el signo de cada factor en $x=-4$, $x=0$ y $x=7$.

\textbf{¿De dónde sale?} Un valor de prueba tomado dentro de cada intervalo del paso 2.

\textbf{¿Por qué?} El signo del producto lo fija la paridad del número de factores negativos (principio del signo de un producto, p.~83).

\begin{center}
\begin{tabular}{|l|c|c|c|}
\hline
Intervalo & $(-\infty,-3)$ & $(-3,6)$ & $(6,\infty)$ \\
\hline
Valor de prueba & $-4$ & $0$ & $7$ \\
Signo de $x-6$ & $-$ & $-$ & $+$ \\
Signo de $x+3$ & $-$ & $+$ & $+$ \\
\hline
Signo de $(x-6)(x+3)$ & $+$ & $-$ & $+$ \\
\hline
\end{tabular}
\end{center}
\end{paso}

\begin{paso}{Paso 4 --- Leer la solución y revisar los extremos}
\textbf{Qué hago:} tomo el intervalo con signo $-$ e incluyo $-3$ y $6$.

\textbf{¿De dónde sale?} Del último renglón de la tabla; los extremos anulan el producto.

\textbf{¿Por qué?} El símbolo es $\le$, así que los valores que dan $0$ sí satisfacen la desigualdad (paso 5 de la guía).
\end{paso}

\begin{center}
\begin{tikzpicture}[scale=0.55]
\draw[->] (-6,0) -- (9,0);
\foreach \x in {-4,-2,0,2,4,6,8} \draw (\x,0.15)--(\x,-0.15) node[below]{\small$\x$};
\draw (-3,0.15)--(-3,-0.15);
\draw[very thick,red] (-3,0) -- (6,0);
\fill[red] (-3,0) circle (3pt); \fill[red] (6,0) circle (3pt);
\end{tikzpicture}
\end{center}

\begin{resp}
$[-3,\,6]$.
\end{resp}

% =====================================================================
\subsection*{Ejercicio 4 --- Cuadrática con coeficiente principal distinto de 1 (libro 43, p. 89)}

\begin{quote}
\textbf{43.} $2x^2+x\ge 1$
\end{quote}

\begin{paso}{Paso 1 --- Mover todo a un lado}
\textbf{Qué hago:} resto $1$ en ambos lados para dejar $0$ a la derecha.

\textbf{¿De dónde sale?} El $1$ es el miembro derecho del enunciado.

\textbf{¿Por qué?} La factorización solo sirve si todos los términos no nulos están de un lado (advertencia de la p.~84).
\[
2x^2+x-1\ge 0.
\]
\end{paso}

\begin{paso}{Paso 2 --- Factorizar}
\textbf{Qué hago:} factorizo por ensayo de factores: $2x^2+x-1=(2x-1)(x+1)$.

\textbf{¿De dónde sale?} $2x\cdot x=2x^2$, $(-1)(1)=-1$, y $2x\cdot 1+(-1)\cdot x=x$.

\textbf{¿Por qué?} En el capítulo~1 los trinomios se factorizan por inspección de los factores del término principal y del independiente.
\[
(2x-1)(x+1)\ge 0.
\]
\end{paso}

\begin{paso}{Paso 3 --- Valores frontera y tabla de signos}
\textbf{Qué hago:} los ceros son $x=\tfrac12$ y $x=-1$; pruebo con $-2$, $0$ y $1$.

\textbf{¿De dónde sale?} De igualar a cero cada factor del paso 2.

\textbf{¿Por qué?} Cada factor mantiene su signo entre ceros consecutivos, así que basta un valor de prueba por intervalo.

\begin{center}
\begin{tabular}{|l|c|c|c|}
\hline
Intervalo & $(-\infty,-1)$ & $\left(-1,\tfrac12\right)$ & $\left(\tfrac12,\infty\right)$ \\
\hline
Valor de prueba & $-2$ & $0$ & $1$ \\
Signo de $2x-1$ & $-$ & $-$ & $+$ \\
Signo de $x+1$ & $-$ & $+$ & $+$ \\
\hline
Signo de $(2x-1)(x+1)$ & $+$ & $-$ & $+$ \\
\hline
\end{tabular}
\end{center}
\end{paso}

\begin{paso}{Paso 4 --- Leer la solución}
\textbf{Qué hago:} tomo los dos intervalos con signo $+$ e incluyo $-1$ y $\tfrac12$.

\textbf{¿De dónde sale?} Del último renglón de la tabla del paso 3.

\textbf{¿Por qué?} El símbolo $\ge$ admite el valor $0$, que se alcanza justo en los ceros de los factores.
\end{paso}

\begin{center}
\begin{tikzpicture}[scale=0.7]
\draw[->] (-4,0) -- (4,0);
\foreach \x in {-3,-2,-1,0,1,2,3} \draw (\x,0.13)--(\x,-0.13) node[below]{\small$\x$};
\draw (0.5,0.13)--(0.5,-0.13);
\node[above] at (0.5,0.15) {\small$\frac12$};
\draw[very thick,red] (-3.6,0) -- (-1,0);
\draw[very thick,red] (0.5,0) -- (3.6,0);
\fill[red] (-1,0) circle (2.6pt); \fill[red] (0.5,0) circle (2.6pt);
\end{tikzpicture}
\end{center}

\begin{resp}
$(-\infty,-1]\cup\left[\dfrac12,\,\infty\right)$.
\end{resp}

% =====================================================================
\subsection*{Ejercicio 5 --- Polinomial con factor repetido (libro 55, p. 89)}

\begin{quote}
\textbf{55.} $(x+3)^2(x-2)(x+5)\ge 0$
\end{quote}

\begin{paso}{Paso 1 --- Valores frontera}
\textbf{Qué hago:} igualo a cero cada factor: $x=-3$ (doble), $x=2$ y $x=-5$.

\textbf{¿De dónde sale?} Del producto ya factorizado del enunciado, con $0$ del otro lado.

\textbf{¿Por qué?} Los ceros de los factores son los únicos puntos donde el producto puede cambiar de signo.
\[
(-\infty,-5),\ (-5,-3),\ (-3,2),\ (2,\infty).
\]
\end{paso}

\begin{paso}{Paso 2 --- Tabla de signos}
\textbf{Qué hago:} evalúo los factores en $x=-6,\,-4,\,0,\,3$. El factor $(x+3)^2$ es $+$ en todos ellos por ser un cuadrado no nulo.

\textbf{¿De dónde sale?} Un valor de prueba por cada intervalo del paso 1.

\textbf{¿Por qué?} El signo del producto lo decide la paridad del número de factores negativos (p.~83); el factor repetido nunca cambia el signo, solo aporta un cero.

\begin{center}
\begin{tabular}{|l|c|c|c|c|}
\hline
Intervalo & $(-\infty,-5)$ & $(-5,-3)$ & $(-3,2)$ & $(2,\infty)$ \\
\hline
Valor de prueba & $-6$ & $-4$ & $0$ & $3$ \\
Signo de $(x+3)^2$ & $+$ & $+$ & $+$ & $+$ \\
Signo de $x-2$ & $-$ & $-$ & $-$ & $+$ \\
Signo de $x+5$ & $-$ & $+$ & $+$ & $+$ \\
\hline
Signo del producto & $+$ & $-$ & $-$ & $+$ \\
\hline
\end{tabular}
\end{center}
\end{paso}

\begin{paso}{Paso 3 --- Revisar los extremos}
\textbf{Qué hago:} incluyo $-5$, $-3$ y $2$ porque el producto vale $0$ en ellos.

\textbf{¿De dónde sale?} Son los ceros hallados en el paso 1.

\textbf{¿Por qué?} El símbolo es $\ge$, y $-3$ queda como punto aislado porque a ambos lados el producto es negativo.
\end{paso}

\begin{center}
\begin{tikzpicture}[scale=0.6]
\draw[->] (-8,0) -- (5,0);
\foreach \x in {-7,-6,-5,-4,-3,-2,-1,0,1,2,3,4} \draw (\x,0.14)--(\x,-0.14);
\node[below] at (-5,-0.14) {\small$-5$};
\node[below] at (-3,-0.14) {\small$-3$};
\node[below] at (0,-0.14) {\small$0$};
\node[below] at (2,-0.14) {\small$2$};
\draw[very thick,red] (-7.6,0) -- (-5,0);
\draw[very thick,red] (2,0) -- (4.6,0);
\fill[red] (-5,0) circle (3pt); \fill[red] (-3,0) circle (3pt); \fill[red] (2,0) circle (3pt);
\end{tikzpicture}
\end{center}

\begin{resp}
$(-\infty,-5]\cup\{-3\}\cup[2,\,\infty)$.
\end{resp}

% =====================================================================
\subsection*{Ejercicio 6 --- Desigualdad cúbica (libro 57, p. 89)}

\begin{quote}
\textbf{57.} $x^3-4x>0$
\end{quote}

\begin{paso}{Paso 1 --- Factorizar}
\textbf{Qué hago:} saco factor común $x$ y luego aplico diferencia de cuadrados a $x^2-4$.

\textbf{¿De dónde sale?} Del lado izquierdo del enunciado, que ya tiene $0$ al otro lado.

\textbf{¿Por qué?} $a^2-b^2=(a-b)(a+b)$ con $a=x$, $b=2$.
\[
x^3-4x=x(x^2-4)=x(x-2)(x+2)>0.
\]
\end{paso}

\begin{paso}{Paso 2 --- Valores frontera y tabla de signos}
\textbf{Qué hago:} los ceros son $-2$, $0$ y $2$; pruebo con $-3,\,-1,\,1,\,3$.

\textbf{¿De dónde sale?} De los tres factores del paso 1.

\textbf{¿Por qué?} Cada factor conserva su signo entre ceros consecutivos.

\begin{center}
\begin{tabular}{|l|c|c|c|c|}
\hline
Intervalo & $(-\infty,-2)$ & $(-2,0)$ & $(0,2)$ & $(2,\infty)$ \\
\hline
Valor de prueba & $-3$ & $-1$ & $1$ & $3$ \\
Signo de $x$ & $-$ & $-$ & $+$ & $+$ \\
Signo de $x-2$ & $-$ & $-$ & $-$ & $+$ \\
Signo de $x+2$ & $-$ & $+$ & $+$ & $+$ \\
\hline
Signo de $x(x-2)(x+2)$ & $-$ & $+$ & $-$ & $+$ \\
\hline
\end{tabular}
\end{center}
\end{paso}

\begin{paso}{Paso 3 --- Leer la solución}
\textbf{Qué hago:} tomo los intervalos con signo $+$ y excluyo los extremos.

\textbf{¿De dónde sale?} Del último renglón del paso 2.

\textbf{¿Por qué?} El símbolo es $>$ estricto, y en $-2,0,2$ el producto vale exactamente $0$.
\end{paso}

\begin{center}
\begin{tikzpicture}[scale=0.7]
\draw[->] (-4,0) -- (4,0);
\foreach \x in {-3,-2,-1,0,1,2,3} \draw (\x,0.13)--(\x,-0.13) node[below]{\small$\x$};
\draw[very thick,red] (-2,0) -- (0,0);
\draw[very thick,red] (2,0) -- (3.6,0);
\draw[red,fill=white,thick] (-2,0) circle (2.6pt);
\draw[red,fill=white,thick] (0,0) circle (2.6pt);
\draw[red,fill=white,thick] (2,0) circle (2.6pt);
\end{tikzpicture}
\end{center}

\begin{resp}
$(-2,\,0)\cup(2,\,\infty)$.
\end{resp}

% =====================================================================
\subsection*{Ejercicio 7 --- Cociente ya factorizado (libro 59, p. 89)}

\begin{quote}
\textbf{59--74 $\blacksquare$ Desigualdades que implican cocientes} \; Resuelva la desigualdad no lineal. Exprese la solución usando notación de intervalos y trace la gráfica del conjunto solución.

\textbf{59.} $\dfrac{x-3}{x+1}\ge 0$
\end{quote}

\begin{paso}{Paso 1 --- Valores frontera del numerador y del denominador}
\textbf{Qué hago:} igualo a cero numerador y denominador: $x=3$ y $x=-1$.

\textbf{¿De dónde sale?} De $x-3$ y $x+1$ del enunciado, que ya está en la forma «cociente $\ge 0$».

\textbf{¿Por qué?} El cociente solo puede cambiar de signo donde se anula el numerador o donde deja de estar definido (paso 3 de la guía, p.~83).
\[
(-\infty,-1),\quad (-1,3),\quad (3,\infty).
\]
\end{paso}

\begin{paso}{Paso 2 --- Tabla de signos}
\textbf{Qué hago:} pruebo con $x=-2$, $x=0$ y $x=4$.

\textbf{¿De dónde sale?} Un valor por cada intervalo del paso 1.

\textbf{¿Por qué?} El signo de un cociente sigue la misma regla de paridad que el de un producto (p.~83).

\begin{center}
\begin{tabular}{|l|c|c|c|}
\hline
Intervalo & $(-\infty,-1)$ & $(-1,3)$ & $(3,\infty)$ \\
\hline
Valor de prueba & $-2$ & $0$ & $4$ \\
Signo de $x-3$ & $-$ & $-$ & $+$ \\
Signo de $x+1$ & $-$ & $+$ & $+$ \\
\hline
Signo de $\dfrac{x-3}{x+1}$ & $+$ & $-$ & $+$ \\
\hline
\end{tabular}
\end{center}
\end{paso}

\begin{paso}{Paso 3 --- Revisar los extremos}
\textbf{Qué hago:} incluyo $x=3$ y excluyo $x=-1$.

\textbf{¿De dónde sale?} $x=3$ anula el numerador; $x=-1$ anula el denominador.

\textbf{¿Por qué?} En $x=3$ el cociente vale $0$, que sí cumple $\ge 0$; en $x=-1$ el cociente no está definido, así que no puede pertenecer a la solución.
\end{paso}

\begin{center}
\begin{tikzpicture}[scale=0.7]
\draw[->] (-4,0) -- (6,0);
\foreach \x in {-3,-2,-1,0,1,2,3,4,5} \draw (\x,0.13)--(\x,-0.13) node[below]{\small$\x$};
\draw[very thick,red] (-3.6,0) -- (-1,0);
\draw[very thick,red] (3,0) -- (5.6,0);
\draw[red,fill=white,thick] (-1,0) circle (2.6pt);
\fill[red] (3,0) circle (2.6pt);
\end{tikzpicture}
\end{center}

\begin{resp}
$(-\infty,-1)\cup[3,\,\infty)$.
\end{resp}

% =====================================================================
\subsection*{Ejercicio 8 --- Cociente con constante al otro lado (libro 61, p. 89)}

\begin{quote}
\textbf{61.} $\dfrac{x}{x+1}>3$
\end{quote}

\begin{paso}{Paso 1 --- Mover todo a un lado y unir en una sola fracción}
\textbf{Qué hago:} resto $3$ y uso $x+1$ como común denominador. \emph{No} multiplico por $x+1$.

\textbf{¿De dónde sale?} El $3$ es el miembro derecho del enunciado; el denominador común es el del enunciado.

\textbf{¿Por qué?} El libro advierte (p.~85) que no se puede multiplicar por $x+1$ sin conocer su signo, porque no se sabría si hay que invertir la desigualdad.
\[
\frac{x}{x+1}-3>0 \;\Longleftrightarrow\; \frac{x-3(x+1)}{x+1}>0 \;\Longleftrightarrow\; \frac{-2x-3}{x+1}>0.
\]
\end{paso}

\begin{paso}{Paso 2 --- Valores frontera}
\textbf{Qué hago:} resuelvo $-2x-3=0$ y $x+1=0$.

\textbf{¿De dónde sale?} Del numerador y del denominador del paso 1.

\textbf{¿Por qué?} Son los únicos puntos donde el cociente puede cambiar de signo.
\[
x=-\tfrac32, \qquad x=-1.
\]
\end{paso}

\begin{paso}{Paso 3 --- Tabla de signos}
\textbf{Qué hago:} pruebo con $x=-2$, $x=-1.2$ y $x=0$.

\textbf{¿De dónde sale?} Un valor dentro de cada uno de los tres intervalos que determinan $-\tfrac32$ y $-1$.

\textbf{¿Por qué?} Cada factor conserva su signo dentro de cada intervalo.

\begin{center}
\begin{tabular}{|l|c|c|c|}
\hline
Intervalo & $\left(-\infty,-\tfrac32\right)$ & $\left(-\tfrac32,-1\right)$ & $(-1,\infty)$ \\
\hline
Valor de prueba & $-2$ & $-1.2$ & $0$ \\
Signo de $-2x-3$ & $+$ & $-$ & $-$ \\
Signo de $x+1$ & $-$ & $-$ & $+$ \\
\hline
Signo de $\dfrac{-2x-3}{x+1}$ & $-$ & $+$ & $-$ \\
\hline
\end{tabular}
\end{center}
\end{paso}

\begin{paso}{Paso 4 --- Leer la solución y verificar}
\textbf{Qué hago:} tomo el intervalo con signo $+$, con extremos abiertos, y compruebo con $x=-1.2$ en el enunciado original.

\textbf{¿De dónde sale?} Del último renglón del paso 3; $-\tfrac32$ anula el numerador y $-1$ el denominador.

\textbf{¿Por qué?} El símbolo es $>$ estricto, así que ni el cero del numerador ni el punto donde el cociente no existe entran.
\[
\frac{-1.2}{-1.2+1}=\frac{-1.2}{-0.2}=6>3 \ \checkmark
\]
\end{paso}

\begin{center}
\begin{tikzpicture}[scale=0.9]
\draw[->] (-3.2,0) -- (1.5,0);
\foreach \x in {-3,-2,-1,0,1} \draw (\x,0.11)--(\x,-0.11) node[below]{\small$\x$};
\draw (-1.5,0.11)--(-1.5,-0.11);
\node[above] at (-1.5,0.13) {\small$-\frac32$};
\draw[very thick,red] (-1.5,0) -- (-1,0);
\draw[red,fill=white,thick] (-1.5,0) circle (2.2pt);
\draw[red,fill=white,thick] (-1,0) circle (2.2pt);
\end{tikzpicture}
\end{center}

\begin{resp}
$\left(-\dfrac32,\,-1\right)$.
\end{resp}

% =====================================================================
\subsection*{Ejercicio 9 --- Cociente con constante al otro lado (libro 63, p. 89)}

\begin{quote}
\textbf{63.} $\dfrac{2x+1}{x-5}\le 3$
\end{quote}

\begin{paso}{Paso 1 --- Pasar todo a un lado}
\textbf{Qué hago:} resto $3$ y combino con denominador común $x-5$.

\textbf{¿De dónde sale?} El $3$ del miembro derecho y el denominador $x-5$ del enunciado.

\textbf{¿Por qué?} No se puede multiplicar por $x-5$ porque su signo es desconocido (p.~85); la guía pide pasar todo a un lado y unir en una fracción.
\[
\frac{2x+1}{x-5}-3\le 0 \;\Longleftrightarrow\; \frac{2x+1-3(x-5)}{x-5}\le 0 \;\Longleftrightarrow\; \frac{-x+16}{x-5}\le 0.
\]
\end{paso}

\begin{paso}{Paso 2 --- Valores frontera y tabla de signos}
\textbf{Qué hago:} los ceros son $x=16$ (numerador) y $x=5$ (denominador); pruebo con $0$, $6$ y $17$.

\textbf{¿De dónde sale?} De $-x+16=0$ y $x-5=0$ en el paso 1.

\textbf{¿Por qué?} El cociente conserva su signo entre esos dos puntos.

\begin{center}
\begin{tabular}{|l|c|c|c|}
\hline
Intervalo & $(-\infty,5)$ & $(5,16)$ & $(16,\infty)$ \\
\hline
Valor de prueba & $0$ & $6$ & $17$ \\
Signo de $-x+16$ & $+$ & $+$ & $-$ \\
Signo de $x-5$ & $-$ & $+$ & $+$ \\
\hline
Signo de $\dfrac{-x+16}{x-5}$ & $-$ & $+$ & $-$ \\
\hline
\end{tabular}
\end{center}
\end{paso}

\begin{paso}{Paso 3 --- Extremos y verificación}
\textbf{Qué hago:} incluyo $x=16$, excluyo $x=5$ y compruebo $x=16$ en el enunciado.

\textbf{¿De dónde sale?} $16$ anula el numerador; $5$ anula el denominador.

\textbf{¿Por qué?} Con $\le$ el valor $0$ sí cumple, pero un punto donde el cociente no está definido nunca pertenece a la solución.
\[
\frac{2(16)+1}{16-5}=\frac{33}{11}=3\le 3 \ \checkmark
\]
\end{paso}

\begin{center}
\begin{tikzpicture}[scale=0.4]
\draw[->] (-6,0) -- (22,0);
\foreach \x in {-4,0,4,8,12,16,20} \draw (\x,0.2)--(\x,-0.2) node[below]{\small$\x$};
\draw (5,0.2)--(5,-0.2);
\node[above] at (5,0.25) {\small$5$};
\draw[very thick,red] (-5.6,0) -- (5,0);
\draw[very thick,red] (16,0) -- (21.6,0);
\draw[red,fill=white,thick] (5,0) circle (4pt);
\fill[red] (16,0) circle (4pt);
\end{tikzpicture}
\end{center}

\begin{resp}
$(-\infty,5)\cup[16,\,\infty)$.
\end{resp}

% =====================================================================
\subsection*{Ejercicio 10 --- Valor absoluto con $\le$ (libro 79, p. 89)}

\begin{quote}
\textbf{75--90 $\blacksquare$ Desigualdades con valor absoluto} \; Resuelva la desigualdad con valor absoluto. Exprese la respuesta usando notación de intervalos y trace la gráfica del conjunto solución.

\textbf{79.} $|\,x-5\,|\le 3$
\end{quote}

\begin{paso}{Paso 1 --- Aplicar la propiedad 2 del valor absoluto}
\textbf{Qué hago:} sustituyo $|A|\le c$ por la doble desigualdad $-c\le A\le c$, con $A=x-5$ y $c=3$.

\textbf{¿De dónde sale?} Del enunciado: la expresión dentro de las barras es $x-5$ y la cota es $3$.

\textbf{¿Por qué?} Propiedad 2 de la p.~86: $|x|\le c \iff -c\le x\le c$ (válida para cualquier expresión algebraica, con $c>0$).
\[
-3\le x-5\le 3.
\]
\end{paso}

\begin{paso}{Paso 2 --- Despejar $x$}
\textbf{Qué hago:} sumo $5$ en los tres miembros.

\textbf{¿De dónde sale?} El $-5$ del miembro central $x-5$.

\textbf{¿Por qué?} Sumar la misma cantidad conserva el sentido en ambas desigualdades (regla 1, p.~82).
\[
2\le x\le 8.
\]
\end{paso}

\begin{center}
\begin{tikzpicture}[scale=0.7]
\draw[->] (-1,0) -- (10,0);
\foreach \x in {0,2,4,6,8} \draw (\x,0.13)--(\x,-0.13) node[below]{\small$\x$};
\draw[very thick,red] (2,0) -- (8,0);
\fill[red] (2,0) circle (2.6pt); \fill[red] (8,0) circle (2.6pt);
\end{tikzpicture}
\end{center}

\begin{resp}
$[2,\,8]$. Geométricamente: los puntos cuya distancia a $5$ no pasa de $3$.
\end{resp}

% =====================================================================
\subsection*{Ejercicio 11 --- Valor absoluto con $\ge$ (libro 83, p. 89)}

\begin{quote}
\textbf{83.} $|\,3x-2\,|\ge 5$
\end{quote}

\begin{paso}{Paso 1 --- Aplicar la propiedad 4 del valor absoluto}
\textbf{Qué hago:} sustituyo $|A|\ge c$ por «$A\le -c$ o $c\le A$», con $A=3x-2$ y $c=5$.

\textbf{¿De dónde sale?} Del enunciado: expresión interior $3x-2$, cota $5$.

\textbf{¿Por qué?} Propiedad 4 de la p.~86: $|x|\ge c \iff x\le -c$ o $c\le x$.
\[
3x-2\le -5 \qquad \text{o} \qquad 3x-2\ge 5.
\]
\end{paso}

\begin{paso}{Paso 2 --- Resolver cada desigualdad}
\textbf{Qué hago:} sumo $2$ y divido entre $3$ en cada rama.

\textbf{¿De dónde sale?} De las dos desigualdades lineales del paso 1.

\textbf{¿Por qué?} $3$ es positivo, así que dividir no invierte ninguno de los dos sentidos (regla 3, p.~82).
\[
3x\le -3 \Rightarrow x\le -1,
\qquad
3x\ge 7 \Rightarrow x\ge \tfrac73.
\]
\end{paso}

\begin{paso}{Paso 3 --- Unir las dos ramas}
\textbf{Qué hago:} escribo la unión de los dos conjuntos.

\textbf{¿De dónde sale?} Del conectivo «o» de la propiedad 4.

\textbf{¿Por qué?} La palabra «o» significa que basta cumplir una de las dos condiciones, lo que corresponde a la unión.
\end{paso}

\begin{center}
\begin{tikzpicture}[scale=0.8]
\draw[->] (-4,0) -- (4.5,0);
\foreach \x in {-3,-2,-1,0,1,2,3,4} \draw (\x,0.12)--(\x,-0.12) node[below]{\small$\x$};
\draw (2.3333,0.12)--(2.3333,-0.12);
\node[above] at (2.3333,0.14) {\small$\frac73$};
\draw[very thick,red] (-3.7,0) -- (-1,0);
\draw[very thick,red] (2.3333,0) -- (4.2,0);
\fill[red] (-1,0) circle (2.4pt); \fill[red] (2.3333,0) circle (2.4pt);
\end{tikzpicture}
\end{center}

\begin{resp}
$(-\infty,-1]\cup\left[\dfrac73,\,\infty\right)$.
\end{resp}

% =====================================================================
\subsection*{Ejercicio 12 --- Valor absoluto que hay que aislar (libro 89, p. 89)}

\begin{quote}
\textbf{89.} $8-|\,2x-1\,|\ge 6$
\end{quote}

\begin{paso}{Paso 1 --- Aislar el valor absoluto}
\textbf{Qué hago:} resto $8$ y luego multiplico por $-1$ invirtiendo el sentido.

\textbf{¿De dónde sale?} El $8$ y el $6$ del enunciado.

\textbf{¿Por qué?} Las propiedades de la p.~86 solo se aplican cuando el valor absoluto está solo de un lado; multiplicar por $-1$ invierte la desigualdad (regla 4, p.~82).
\[
-|2x-1|\ge -2 \;\Longleftrightarrow\; |2x-1|\le 2.
\]
\end{paso}

\begin{paso}{Paso 2 --- Aplicar la propiedad 2}
\textbf{Qué hago:} convierto $|2x-1|\le 2$ en la doble desigualdad $-2\le 2x-1\le 2$.

\textbf{¿De dónde sale?} Del resultado del paso 1, con $A=2x-1$ y $c=2$.

\textbf{¿Por qué?} Propiedad 2 de la p.~86.
\end{paso}

\begin{paso}{Paso 3 --- Despejar $x$}
\textbf{Qué hago:} sumo $1$ y divido entre $2$ en los tres miembros.

\textbf{¿De dónde sale?} Del $-1$ y del coeficiente $2$ del miembro central.

\textbf{¿Por qué?} Dividir entre el número positivo $2$ conserva ambos sentidos.
\[
-1\le 2x\le 3 \;\Longleftrightarrow\; -\tfrac12\le x\le \tfrac32.
\]
\end{paso}

\begin{center}
\begin{tikzpicture}[scale=1.1]
\draw[->] (-1.6,0) -- (2.4,0);
\foreach \x in {-1,0,1,2} \draw (\x,0.09)--(\x,-0.09) node[below]{\small$\x$};
\draw (-0.5,0.09)--(-0.5,-0.09); \draw (1.5,0.09)--(1.5,-0.09);
\node[above] at (-0.5,0.11) {\small$-\frac12$};
\node[above] at (1.5,0.11) {\small$\frac32$};
\draw[very thick,red] (-0.5,0) -- (1.5,0);
\fill[red] (-0.5,0) circle (1.9pt); \fill[red] (1.5,0) circle (1.9pt);
\end{tikzpicture}
\end{center}

\begin{resp}
$\left[-\dfrac12,\,\dfrac32\right]$.
\end{resp}

% =====================================================================
\subsection*{Ejercicio 13 --- De la frase a la desigualdad (libro 93, p. 89)}

\begin{quote}
\textbf{91--94 $\blacksquare$ Desigualdades con valor absoluto} \; Se proporciona una frase que describe un conjunto de números reales. Exprese la frase como una desigualdad que contenga un valor absoluto.

\textbf{93.} Todos los números reales $x$ menores de 5 unidades desde 7
\end{quote}

\begin{paso}{Paso 1 --- Traducir «distancia» a valor absoluto}
\textbf{Qué hago:} escribo la distancia entre $x$ y $7$ como $|x-7|$ y le impongo la cota estricta $5$.

\textbf{¿De dónde sale?} «desde 7» fija el centro; «menores de 5 unidades» fija la cota y el símbolo $<$.

\textbf{¿Por qué?} La distancia entre dos puntos $a$ y $b$ de la recta real es $|b-a|$ (definición del cap.~1, usada en la figura~8, p.~86).
\[
|x-7|<5.
\]
\end{paso}

\begin{paso}{Paso 2 --- Comprobar resolviendo}
\textbf{Qué hago:} aplico la propiedad 1 y verifico que el conjunto obtenido es el descrito.

\textbf{¿De dónde sale?} De $|x-7|<5$ con $c=5$.

\textbf{¿Por qué?} Propiedad 1 de la p.~86: $|x|<c \iff -c<x<c$.
\[
-5<x-7<5 \;\Longleftrightarrow\; 2<x<12,
\]
que son, en efecto, los puntos a menos de $5$ unidades de $7$.
\end{paso}

\begin{resp}
$|x-7|<5$, cuyo conjunto solución es $(2,\,12)$.
\end{resp}

% =====================================================================
\subsection*{Ejercicio 14 --- Dominio de una raíz (libro 101, p. 89)}

\begin{quote}
\textbf{101--104 $\blacksquare$ Dominio} \; Determine los valores de la variable para la cual la expresión está definida como número real.

\textbf{101.} $\sqrt{x^2-9}$
\end{quote}

\begin{paso}{Paso 1 --- Plantear la restricción del radicando}
\textbf{Qué hago:} exijo $x^2-9\ge 0$.

\textbf{¿De dónde sale?} $x^2-9$ es el radicando de la raíz cuadrada del enunciado.

\textbf{¿Por qué?} Una raíz de índice par solo está definida en los reales para argumentos no negativos.
\end{paso}

\begin{paso}{Paso 2 --- Factorizar y armar la tabla de signos}
\textbf{Qué hago:} factorizo por diferencia de cuadrados y pruebo con $-4$, $0$ y $4$.

\textbf{¿De dónde sale?} $x^2-9=(x-3)(x+3)$ con $a=x$, $b=3$; los ceros $\pm 3$ dan los tres intervalos.

\textbf{¿Por qué?} Es la guía de desigualdades no lineales (p.~83) aplicada a $(x-3)(x+3)\ge 0$.

\begin{center}
\begin{tabular}{|l|c|c|c|}
\hline
Intervalo & $(-\infty,-3)$ & $(-3,3)$ & $(3,\infty)$ \\
\hline
Valor de prueba & $-4$ & $0$ & $4$ \\
Signo de $x-3$ & $-$ & $-$ & $+$ \\
Signo de $x+3$ & $-$ & $+$ & $+$ \\
\hline
Signo de $(x-3)(x+3)$ & $+$ & $-$ & $+$ \\
\hline
\end{tabular}
\end{center}
\end{paso}

\begin{paso}{Paso 3 --- Leer el dominio}
\textbf{Qué hago:} tomo los intervalos con signo $+$ e incluyo $\pm 3$.

\textbf{¿De dónde sale?} Del último renglón del paso 2.

\textbf{¿Por qué?} En $x=\pm3$ el radicando vale $0$ y $\sqrt{0}=0$ sí es un número real, así que el símbolo $\ge$ los admite.
\end{paso}

\begin{center}
\begin{tikzpicture}[scale=0.7]
\draw[->] (-5,0) -- (5,0);
\foreach \x in {-4,-3,-2,-1,0,1,2,3,4} \draw (\x,0.13)--(\x,-0.13) node[below]{\small$\x$};
\draw[very thick,red] (-4.6,0) -- (-3,0);
\draw[very thick,red] (3,0) -- (4.6,0);
\fill[red] (-3,0) circle (2.6pt); \fill[red] (3,0) circle (2.6pt);
\end{tikzpicture}
\end{center}

\begin{resp}
$(-\infty,-3]\cup[3,\,\infty)$.
\end{resp}

% =====================================================================
\subsection*{Ejercicio 15 --- Aplicación: comparar dos planes de renta (libro 111, p. 90)}

\begin{quote}
\textbf{111. Costo de la renta de un auto} \; Una compañía de renta de autos ofrece dos planes para la renta de un auto.

Plan A:\quad 30 dólares por día y 10 centavos por milla

Plan B:\quad 50 dólares por día con kilometraje ilimitado

¿Para qué intervalo de millas es más económico el plan B?
\end{quote}

\begin{paso}{Paso 1 --- Identificar la variable}
\textbf{Qué hago:} defino $x=$ número de millas recorridas en un día.

\textbf{¿De dónde sale?} La pregunta es «para qué intervalo de millas», así que la incógnita son las millas.

\textbf{¿Por qué?} Es el primer paso del modelado con desigualdades del libro (ejemplo 8, p.~87).
\end{paso}

\begin{paso}{Paso 2 --- Traducir a álgebra}
\textbf{Qué hago:} escribo el costo diario de cada plan.

\textbf{¿De dónde sale?} Plan A: cuota fija $30$ más $0.10$ dólares por cada una de las $x$ millas. Plan B: $50$ dólares fijos, sin cargo por milla.

\textbf{¿Por qué?} 10 centavos $=0.10$ dólares, y el kilometraje ilimitado significa que el costo no depende de $x$.
\[
\text{costo A}=30+0.10x, \qquad \text{costo B}=50.
\]
\end{paso}

\begin{paso}{Paso 3 --- Formular el modelo}
\textbf{Qué hago:} impongo que el plan B sea más barato que el A.

\textbf{¿De dónde sale?} De los costos del paso 2 y de la frase «es más económico el plan B».

\textbf{¿Por qué?} «Más económico» significa costo estrictamente menor.
\[
50<30+0.10x.
\]
\end{paso}

\begin{paso}{Paso 4 --- Resolver}
\textbf{Qué hago:} resto $30$ y divido entre $0.10$.

\textbf{¿De dónde sale?} De la desigualdad del paso 3.

\textbf{¿Por qué?} $0.10$ es positivo, así que dividir conserva el sentido (regla 3, p.~82).
\[
20<0.10x \;\Longleftrightarrow\; 200<x.
\]
\end{paso}

\begin{paso}{Paso 5 --- Verificar}
\textbf{Qué hago:} comparo los costos en $x=300$ y en $x=100$.

\textbf{¿De dónde sale?} Un valor dentro y otro fuera del intervalo del paso 4.

\textbf{¿Por qué?} Confirma que la desigualdad se planteó en el sentido correcto.
\[
x=300:\ 30+0.10(300)=60>50 \ \checkmark
\qquad
x=100:\ 30+0.10(100)=40<50 \ \checkmark
\]
\end{paso}

\begin{resp}
El plan B resulta más económico si se recorren más de $200$ millas al día: $x\in(200,\infty)$. (En $x=200$ ambos planes cuestan $50$ dólares.)
\end{resp}

\end{document}
